% généré depuis source/10_perspectives.md — ne pas modifier à la main
\chapter{Perspectives}\label{sec:perspectives}

\chapintro

Ce chapitre présente les développements envisagés pour les Versions 2 et 3. \textbf{Aucun d'entre eux n'est réalisé à ce stade.} Pour chacun, les huit points suivants sont précisés : le développement envisagé, sa motivation, la limite actuelle qui le justifie, les données et ressources nécessaires, les méthodes envisageables, les expérimentations de validation, les bénéfices attendus et les difficultés prévisibles.

\section{Modèles de dégradation probabilistes plus avancés}\label{sec:p-degradation}

\begin{itemize}
\item \textbf{Développement :} remplacer la base polynomiale par un modèle exponentiel à coefficients aléatoires, puis par un processus de Wiener à dérive non linéaire \cite{gebraeel2005,si2012}.
\item \textbf{Motivation :} ces modèles décrivent mieux une dégradation qui s'accélère en fin de vie et intègrent naturellement le bruit de dégradation.
\item \textbf{Limite actuelle :} le polynôme de degré 2 extrapole mal la fin de vie, ce qui contribue au biais de +13,5 cycles.
\item \textbf{Données et ressources :} les mêmes données FD001 suffisent ; aucune ressource matérielle particulière.
\item \textbf{Méthodes :} régression bayésienne sur le logarithme de l'indicateur, estimation par maximum de vraisemblance des paramètres du processus de Wiener, loi du premier temps de passage.
\item \textbf{Validation :} comparaison du RMSE, du biais et du score NASA avec le modèle actuel, sur les mêmes moteurs de test.
\item \textbf{Bénéfices attendus :} une réduction du biais et un modèle plus fidèle à la physique de l'usure.
\item \textbf{Difficultés :} l'indicateur de santé peut devenir négatif, ce qui complique le passage au logarithme ; un recalage de l'indicateur sera nécessaire.
\end{itemize}

\section{Amélioration de l'estimation du RUL et de la quantification des incertitudes}\label{sec:p-incertitude}

\begin{itemize}
\item \textbf{Développement :} propager le bruit de mesure dans les trajectoires simulées et recalibrer les intervalles de prédiction.
\item \textbf{Motivation :} un intervalle n'est utile pour planifier la maintenance que s'il couvre réellement le RUL dans la proportion annoncée.
\item \textbf{Limite actuelle :} l'intervalle P10–P90 ne couvre le RUL réel que dans 54 \% des cas.
\item \textbf{Données et ressources :} les moteurs d'entraînement, dont une partie serait réservée à la calibration.
\item \textbf{Méthodes :} simulation du bruit dans les trajectoires, recalibration par quantiles sur un jeu de validation, prédiction conforme.
\item \textbf{Validation :} courbes de calibration à plusieurs niveaux, sur des moteurs jamais utilisés pour l'apprentissage ni pour la calibration.
\item \textbf{Bénéfices attendus :} des intervalles fiables, utilisables pour un remplacement fondé sur un niveau de risque maîtrisé.
\item \textbf{Difficultés :} des intervalles mieux calibrés seront aussi plus larges ; il faudra trouver un compromis entre couverture et utilité.
\end{itemize}

\section{Étude approfondie des indicateurs B10 et B50}\label{sec:p-bp}

\begin{itemize}
\item \textbf{Développement :} calculer des intervalles de confiance des durées $B_{10}$ et $B_{50}$ et étudier leur sensibilité au choix de la loi.
\item \textbf{Motivation :} ces durées servent à fixer des intervalles de maintenance ; leur incertitude doit être connue.
\item \textbf{Limite actuelle :} les durées $B_p$ sont données sans intervalle de confiance, et deux lois équivalentes (Weibull 3P et log-normale) peuvent en donner des valeurs différentes.
\item \textbf{Données et ressources :} les 200 durées de vie, censures comprises.
\item \textbf{Méthodes :} bootstrap non paramétrique \cite{efron1979}, méthode delta, ajustement de la Weibull 3P avec censure.
\item \textbf{Validation :} stabilité des intervalles selon le nombre de rééchantillonnages ; comparaison des valeurs selon la loi retenue.
\item \textbf{Bénéfices attendus :} des recommandations de maintenance accompagnées de leur incertitude.
\item \textbf{Difficultés :} l'estimation du paramètre de position par maximum de vraisemblance est délicate, en particulier sur des échantillons rééchantillonnés.
\end{itemize}

\section{Intégration de covariables externes}\label{sec:p-covariables}

\begin{itemize}
\item \textbf{Développement :} tenir compte des conditions de vol dans la fiabilité et dans le pronostic.
\item \textbf{Motivation :} en exploitation réelle, la durée de vie dépend fortement des conditions d'utilisation.
\item \textbf{Limite actuelle :} FD001 ne comporte qu'une seule condition opératoire ; cette perspective nécessite les sous-ensembles \textbf{FD002 et FD004}, qui en comportent six.
\item \textbf{Données et ressources :} les sous-ensembles FD002 et FD004, plus volumineux, et un temps de calcul accru.
\item \textbf{Méthodes :} modèle à risques proportionnels de Cox \cite{cox1972}, normalisation des capteurs par condition opératoire.
\item \textbf{Validation :} comparaison des performances avec et sans covariables, sur les jeux de test officiels.
\item \textbf{Bénéfices attendus :} des modèles transposables à des conditions d'exploitation variées.
\item \textbf{Difficultés :} la normalisation par condition opératoire et la présence de deux modes de défaillance dans FD004.
\end{itemize}

\section{Explicabilité des modèles d'apprentissage}\label{sec:p-xai}

\begin{itemize}
\item \textbf{Développement :} expliquer les prédictions du classifieur et du LSTM.
\item \textbf{Motivation :} une décision de maintenance doit pouvoir être justifiée auprès des équipes techniques et des autorités.
\item \textbf{Limite actuelle :} le LSTM et le gradient boosting sont des modèles opaques ; seuls les poids de l'indicateur de santé sont interprétables.
\item \textbf{Données et ressources :} les modèles existants et les données de test.
\item \textbf{Méthodes :} valeurs de Shapley, méthode SHAP \cite{lundberg2017}, et attributions temporelles pour le LSTM.
\item \textbf{Validation :} cohérence des capteurs les plus influents avec la physique de la dégradation du compresseur haute pression.
\item \textbf{Bénéfices attendus :} une confiance accrue des utilisateurs et une vérification de l'absence de corrélations parasites.
\item \textbf{Difficultés :} le coût de calcul des explications sur des séquences et l'interprétation d'attributions corrélées entre capteurs.
\end{itemize}

\section{Stratégies de pondération}\label{sec:p-ponderation}

\begin{itemize}
\item \textbf{Développement :} améliorer la fusion des capteurs dans l'indicateur de santé et la pondération des deux modèles de RUL.
\item \textbf{Motivation :} une moyenne simple des deux modèles n'améliore pas le LSTM seul.
\item \textbf{Limite actuelle :} l'indicateur de santé est une régression linéaire simple, et la combinaison des modèles est une moyenne non pondérée (RMSE 16,91 contre 15,35 pour le LSTM seul).
\item \textbf{Données et ressources :} les moteurs d'entraînement, avec un jeu de validation dédié à l'apprentissage des poids.
\item \textbf{Méthodes :} sélection par monotonie et pronosticabilité \cite{coble2009}, pondération selon l'incertitude, empilement de modèles.
\item \textbf{Validation :} validation croisée par groupes de moteurs, comme pour la classification.
\item \textbf{Bénéfices attendus :} un indicateur plus monotone et une combinaison au moins aussi précise que le meilleur modèle.
\item \textbf{Difficultés :} le risque de sur-apprentissage des poids sur un nombre limité de moteurs.
\end{itemize}

\section{Analyses de sensibilité et de robustesse}\label{sec:p-sensibilite}

\begin{itemize}
\item \textbf{Développement :} mesurer l'effet des choix de paramètres sur les résultats.
\item \textbf{Motivation :} plusieurs paramètres sont fixés par convention et non optimisés.
\item \textbf{Limite actuelle :} le seuil de criticité (30 cycles), la taille des fenêtres, le plafond du LSTM et les hypothèses de réparation et de coût sont fixés sans étude de sensibilité.
\item \textbf{Données et ressources :} les données actuelles ; un temps de calcul proportionnel au nombre de configurations testées.
\item \textbf{Méthodes :} plans d'expériences, variation d'un facteur à la fois, analyse de sensibilité globale.
\item \textbf{Validation :} stabilité des conclusions (choix de la loi, âge optimal, classement des modèles) selon les paramètres.
\item \textbf{Bénéfices attendus :} des conclusions dont la robustesse est démontrée.
\item \textbf{Difficultés :} le nombre de combinaisons croît vite, en particulier si le LSTM doit être réentraîné.
\end{itemize}

\section{Benchmarks élargis, apprentissage profond et Transformers}\label{sec:p-benchmark}

\begin{itemize}
\item \textbf{Développement :} comparer davantage de familles de modèles : réseaux convolutifs, GRU, Transformers et forêts de survie aléatoires.
\item \textbf{Motivation :} situer les deux modèles actuels par rapport à l'état de l'art \cite{babu2016,li2018,ramasso2014}.
\item \textbf{Limite actuelle :} seules deux approches sont comparées.
\item \textbf{Données et ressources :} les quatre sous-ensembles C-MAPSS, et éventuellement un processeur graphique.
\item \textbf{Méthodes :} réseaux convolutifs temporels, GRU, Transformers \cite{vaswani2017}, forêts de survie aléatoires \cite{ishwaran2008}, avec un protocole d'évaluation commun.
\item \textbf{Validation :} mêmes métriques, mêmes moteurs de test et plusieurs graines pour les modèles stochastiques.
\item \textbf{Bénéfices attendus :} un choix de modèle argumenté par une comparaison équitable.
\item \textbf{Difficultés :} le coût de calcul, et le risque d'optimiser les modèles sur le jeu de test ; les Transformers ne seront retenus que si leur gain le justifie au regard de leur complexité.
\end{itemize}

\section{Passage à des données industrielles réelles}\label{sec:p-donnees}

\begin{itemize}
\item \textbf{Développement :} appliquer la chaîne à des données de capteurs réelles.
\item \textbf{Motivation :} C-MAPSS est une simulation ; les données réelles sont plus bruitées, incomplètes et comportent des maintenances intermédiaires.
\item \textbf{Limite actuelle :} toutes les données utilisées sont simulées, et les temps de réparation le sont aussi.
\item \textbf{Données et ressources :} un partenariat industriel, des historiques de capteurs et de maintenance, et un cadre de confidentialité.
\item \textbf{Méthodes :} nettoyage et synchronisation des données, traitement des valeurs manquantes, prise en compte des maintenances imparfaites.
\item \textbf{Validation :} comparaison des prédictions avec les défaillances et les déposes réellement constatées.
\item \textbf{Bénéfices attendus :} une validation de l'approche dans un contexte réel.
\item \textbf{Difficultés :} l'accès aux données, leur qualité et le faible nombre de défaillances observées.
\end{itemize}

\section{Modèles hybrides et apprentissage informé par la physique}\label{sec:p-hybride}

\begin{itemize}
\item \textbf{Développement :} combiner un modèle physique de la dégradation et un modèle appris.
\item \textbf{Motivation :} les modèles purement statistiques n'intègrent pas les lois physiques connues du moteur.
\item \textbf{Limite actuelle :} les modèles actuels sont entièrement guidés par les données.
\item \textbf{Données et ressources :} une connaissance des lois de dégradation du compresseur et, idéalement, un modèle thermodynamique simplifié du moteur.
\item \textbf{Méthodes :} réseaux informés par la physique \cite{raissi2019}, modèles hybrides où le réseau corrige un modèle physique.
\item \textbf{Validation :} précision, mais aussi cohérence physique des prédictions, par exemple la monotonie de la dégradation.
\item \textbf{Bénéfices attendus :} une meilleure extrapolation avec peu de données et des prédictions physiquement plausibles.
\item \textbf{Difficultés :} la formulation des contraintes physiques et l'équilibre entre les termes de la fonction de coût.
\end{itemize}

\section{Optimisation avancée de la planification de la maintenance}\label{sec:p-planification}

\begin{itemize}
\item \textbf{Développement :} passer de la politique d'âge fixe à une planification dynamique fondée sur le RUL prédit.
\item \textbf{Motivation :} la politique actuelle sacrifie en moyenne 75,6 cycles de vie utile par moteur.
\item \textbf{Limite actuelle :} la politique optimisée est identique pour tous les moteurs et ignore leur état.
\item \textbf{Données et ressources :} les distributions du RUL, des contraintes opérationnelles (disponibilité des équipes, des pièces et des créneaux).
\item \textbf{Méthodes :} politiques de remplacement conditionnelles au risque, programmation linéaire en nombres entiers, optimisation sous incertitude.
\item \textbf{Validation :} simulation de la flotte sur un horizon long, comparaison des coûts et du nombre de défaillances en service.
\item \textbf{Bénéfices attendus :} une réduction du coût et de la vie sacrifiée par rapport à la politique d'âge fixe.
\item \textbf{Difficultés :} la nécessité d'intervalles de RUL bien calibrés, d'où le lien avec la section~\ref{sec:p-incertitude}.
\end{itemize}

\section{Jumeau numérique}\label{sec:p-jumeau}

\begin{itemize}
\item \textbf{Développement :} intégrer les modèles dans un jumeau numérique de la flotte, mis à jour en continu.
\item \textbf{Motivation :} suivre l'état de chaque moteur en temps réel et simuler des scénarios de maintenance \cite{tao2019}.
\item \textbf{Limite actuelle :} les modèles sont appliqués hors ligne, sur des fichiers.
\item \textbf{Données et ressources :} un flux de données en continu et une infrastructure logicielle de déploiement.
\item \textbf{Méthodes :} mise à jour bayésienne en ligne (déjà compatible avec le modèle de dégradation), architecture orientée événements.
\item \textbf{Validation :} latence de mise à jour et stabilité des prédictions dans le temps.
\item \textbf{Bénéfices attendus :} une aide à la décision en temps réel.
\item \textbf{Difficultés :} l'intégration aux systèmes existants et la maintenance du logiciel.
\end{itemize}

\section{Étude économique industrielle complète}\label{sec:p-economie}

\begin{itemize}
\item \textbf{Développement :} remplacer les coûts relatifs supposés par des coûts réels.
\item \textbf{Motivation :} l'âge optimal et le gain dépendent directement du rapport des coûts.
\item \textbf{Limite actuelle :} le rapport $C_c / C_p$ est une hypothèse d'illustration, et l'immobilisation n'est pas chiffrée.
\item \textbf{Données et ressources :} des coûts de pièces, de main-d'œuvre et d'immobilisation, auprès d'un exploitant.
\item \textbf{Méthodes :} modèle de coût complet incluant l'immobilisation, les stocks et les pénalités, analyse de sensibilité aux coûts.
\item \textbf{Validation :} comparaison avec les coûts historiques d'un exploitant.
\item \textbf{Bénéfices attendus :} une recommandation économiquement justifiée.
\item \textbf{Difficultés :} la confidentialité et la variabilité des coûts.
\end{itemize}

\section{Déploiement et validation en environnement industriel}\label{sec:p-deploiement}

\begin{itemize}
\item \textbf{Développement :} déployer le tableau de bord et l'API auprès d'utilisateurs réels.
\item \textbf{Motivation :} vérifier l'utilité de l'outil dans le processus de décision.
\item \textbf{Limite actuelle :} l'application fonctionne en local, avec des données figées.
\item \textbf{Données et ressources :} un serveur, une authentification des utilisateurs et une connexion aux données réelles.
\item \textbf{Méthodes :} conteneurisation, intégration continue s'appuyant sur la suite de tests existante, surveillance de la dérive des données.
\item \textbf{Validation :} retours des utilisateurs et indicateurs de maintenance avant et après le déploiement.
\item \textbf{Bénéfices attendus :} un outil réellement utilisé pour la maintenance.
\item \textbf{Difficultés :} la sécurité, la fiabilité du service et la conduite du changement.
\end{itemize}

\chapconclusion

Ces perspectives vont des plus directes, qui s'appuient sur les données actuelles, aux plus ambitieuses, qui demandent des données industrielles ou un déploiement réel. Les premières alimenteront la Version 2 du rapport, en fonction des remarques de l'encadrant.
